\documentclass[12pt]{article}
\usepackage[utf8]{inputenc}
\usepackage[T1]{fontenc}
\usepackage{hyperref}
\title{Semaq Beri Language Resources for K-12 Curriculum Development}
\author{Compiled by Claude Sonnet 5.5 (claude-sonnet-5-5) for Andr\'as Kornai}
\date{}
\begin{document}
\maketitle

\section{Introduction}
Semaq Beri (ISO 639-3 szc; also Semoq Beri, Semaq Bri) is a Southern Aslian (Austroasiatic) language of the Malay Peninsula, spoken in Pahang and Terengganu, Malaysia. It belongs to the same Southern Aslian group as Semelai, Temoq and Mah Meri. The Endangered Languages Project lists it as threatened, with about 2,500 speakers. Almost all documentation comes from one researcher, Nicole Kruspe.

\section{Primary Educational Resources}
No primer, textbook or school programme for Semaq Beri was found in search results.

\section{Linguistic and Reference Materials}
Kruspe's grammar sketch ``Semaq Beri'' appears in \emph{The Handbook of Austroasiatic Languages} (Brill, 2015; Wikipedia gives 2014 for a preliminary description): \url{https://glottolog.org/resource/reference/id/450297}. I found no published dictionary. The Kruspe Collection at the Humlab archive in Lund holds a Toolbox-format lexicon for the Sungai Berua community, Terengganu (2011): \url{https://archive.humlab.lu.se/flat/islandora/object/lat%3A10050_00_0000_0000_0004_09D7_6}. Overview: \url{https://en.wikipedia.org/wiki/Semaq_Beri_language}.

\section{Religious and Community Resources}
No scripture or community-organisation resources were found in search.

\section{Digital and Audio Resources}
The documentation project ``Semoq Beri and Batek'' (HGLC, alongside the DoBeS project Tongues of the Semang) has pages with project links: \url{https://dobes.mpi.nl/projects/semaq_beri_batek/links/}. Further Kruspe Collection items, including Batek lexicons with Semaq Beri comparisons, are in the Humlab archive. I did not verify what audio or video the archive holds or whether access is open. Catalogue entry: \url{https://www.endangeredlanguages.com/es/node/92472}.

\section{Curriculum Development Context}
The language is reported as still being used as a first language in the community, but little pedagogical material exists. The existing documentation (grammar sketch, lexicon, recordings) would be a starting point for teacher-and-community-led materials. Access conditions for the archives would need to be checked.
\end{document}